\begin{center}
\Large{概要}
\end{center}

聴音トレーニングとは，聴覚を用いて音程やリズム等の音楽的要素を識別するトレーニングで，こうしたトレーニングによる音程認識能力の習得は音楽教育の基盤である．これまで，本訓練は指導者に依存していたが，近年ではソフトウェアによる独習へと移行してきており，Virtual Reality（VR）技術を用いた多感覚的アプローチが注目されている．特に，VR特有の空間音響は音程認識の新たな手がかりとなり，従来の学習環境よりも高い学習効果をもたらす可能性が示唆されている．さらに、既存の平面的なディスプレイを用いた学習では，視覚と聴覚の連携は限定的であり，実空間での音源定位や身体性を伴う知覚体験は提供されていない。そこで本研究では、人間が持つ音高の変化を空間的な「高さ」や「配置」として認知する特性を活用し、VR 空間において音源を立体的に配置することで，この空間認知能力を聴音学習に直接的に利用するアプローチに注目した。これにより、単なる知識の習得を超え，音程を直感的かつ身体的に捉える新たな認知スキームの獲得につながると考えられる．


そこで，本研究では，VR空間内で音源を立体的に配置することで学習者に与える影響を明らかにすることを目的とし，VR技術を適用した聴音トレーニングの音源配置の違いについて検討を行った．具体的には，既存手法である平面的インターフェースに対応した前方半円状に音源を配置するシステムと，新たに提案する音の高さを空間的距離や方位に拡張させた円状配置システムといった2つのシステム間で，NASA-TLXによる作業負荷やテスト時間，正答率を比較検証した．
